\documentclass[10pt,a4paper]{amsart}
\usepackage{amsmath,amssymb,amsthm,graphicx}
\usepackage{verbatim}
\newenvironment{reference}{\comment}{\endcomment}
%\newenvironment{reference}{}{}
\newenvironment{slogan}{\comment}{\endcomment}
\newenvironment{history}{\comment}{\endcomment}

% For commutative diagrams we use Xy-pic
\usepackage[all]{xy}

% We use 2cell for 2-commutative diagrams.
\xyoption{2cell}
% Theorem environments.
%
\theoremstyle{plain}
\newtheorem{theorem}[subsection]{Theorem}
\newtheorem{proposition}[subsection]{Proposition}
\newtheorem{lemma}[subsection]{Lemma}

\theoremstyle{definition}
\newtheorem{definition}[subsection]{Definition}
\newtheorem{example}[subsection]{Example}
\newtheorem{exercise}[subsection]{Exercise}
\newtheorem{situation}[subsection]{Situation}

\theoremstyle{remark}
\newtheorem{remark}[subsection]{Remark}
\newtheorem{remarks}[subsection]{Remarks}

\numberwithin{equation}{subsection}

% Macros
%
\def\lim{\mathop{\mathrm{lim}}\nolimits}
\def\colim{\mathop{\mathrm{colim}}\nolimits}
\def\Spec{\mathop{\mathrm{Spec}}}
\def\Hom{\mathop{\mathrm{Hom}}\nolimits}
\def\Ext{\mathop{\mathrm{Ext}}\nolimits}
\def\SheafHom{\mathop{\mathcal{H}\!\mathit{om}}\nolimits}
\def\SheafExt{\mathop{\mathcal{E}\!\mathit{xt}}\nolimits}
\def\Sch{\mathit{Sch}}
\def\Mor{\mathop{\mathrm{Mor}}\nolimits}
\def\Ob{\mathop{\mathrm{Ob}}\nolimits}
\def\Sh{\mathop{\mathit{Sh}}\nolimits}
\def\NL{\mathop{N\!L}\nolimits}
\def\CH{\mathop{\mathrm{CH}}\nolimits}
\def\proetale{{pro\text{-}\acute{e}tale}}
\def\etale{{\acute{e}tale}}
\def\QCoh{\mathit{QCoh}}
\def\Ker{\mathop{\mathrm{Ker}}}
\def\Im{\mathop{\mathrm{Im}}}
\def\Coker{\mathop{\mathrm{Coker}}}
\def\Coim{\mathop{\mathrm{Coim}}}

% Boxtimes
%
\DeclareMathSymbol{\boxtimes}{\mathbin}{AMSa}{"02}

%
% Macros for moduli stacks/spaces
%
\def\QCohstack{\mathcal{QC}\!\mathit{oh}}
\def\Cohstack{\mathcal{C}\!\mathit{oh}}
\def\Spacesstack{\mathcal{S}\!\mathit{paces}}
\def\Quotfunctor{\mathrm{Quot}}
\def\Hilbfunctor{\mathrm{Hilb}}
\def\Curvesstack{\mathcal{C}\!\mathit{urves}}
\def\Polarizedstack{\mathcal{P}\!\mathit{olarized}}
\def\Complexesstack{\mathcal{C}\!\mathit{omplexes}}
% \Pic is the operator that assigns to X its picard group, usage \Pic(X)
% \Picardstack_{X/B} denotes the Picard stack of X over B
% \Picardfunctor_{X/B} denotes the Picard functor of X over B
\def\Pic{\mathop{\mathrm{Pic}}\nolimits}
\def\Picardstack{\mathcal{P}\!\mathit{ic}}
\def\Picardfunctor{\mathrm{Pic}}
\def\Deformationcategory{\mathcal{D}\!\mathit{ef}}
\usepackage{cleveref}
\crefname{originalreference}{original reference}{original references}
\Crefname{originalreference}{Original reference}{Original references}
\makeatletter
\newlabel{sites-lemma-localize-morphism-topoi}{{04H1}{source}{}{sourcepointer-1}{}}
\newlabel{sites-lemma-localize-morphism-topoi@cref}{{[originalreference][0][]04H1}{[1][0][]source}{}{}{}}
\makeatother
\crefname{originallemma}{original lemma}{original lemmas}
\Crefname{originallemma}{Original lemma}{Original lemmas}
\makeatletter
\newlabel{lemma-localize-ringed-topos}{{04IE}{source}{}{sourcepointer-2}{}}
\newlabel{lemma-localize-ringed-topos@cref}{{[originallemma][0][]04IE}{[1][0][]source}{}{}{}}
\makeatother
\title{The relocalization square of ringed topoi: a morphism of ringed topoi relocalizes along a morphism of sheaves}
\author{The Stacks Project Authors}
\date{stacks-project:sites-modules.tex:local-snapshot-20261012; GNU FDL 1.2 or later}
\begin{document}
\maketitle
\noindent\textit{Editorial note.} All delivered mathematical text is the author's own, copied verbatim from the pinned Stacks Project chapter file sites-modules.tex; no statement or proof step was rewritten, reordered or removed.\par\medskip
\noindent\textit{Editorial note.} The counted span cites two external items (the ringed-topos relocalization lemma and a section), recorded in assembly.external\_references with canonical tags as navigation pointers.\par\medskip
\noindent\textit{Allowed standard prerequisites (editorial).}\begin{itemize}
\item Localization of a morphism of topoi along a morphism of sheaves. Stacks Project, Sites, Lemma sites-lemma-localize-morphism-topoi (tag 00YW). Supplies the diagram of underlying topoi, the identity f'\_* j\_F\textasciicircum{}\{-1\} = j\_G\textasciicircum{}\{-1\} f\_* and the description of (f')\textasciicircum{}\{-1\}.
\item Definition of a morphism of ringed topoi. Stacks Project, Modules on Sites, section on morphisms of ringed topoi (structure map of a morphism). Gives the data that must be verified to conclude that the relocalized square is a square of ringed topoi.
\end{itemize}\medskip
\ifdefined\hypertarget\hypertarget{sourcepointer-1}{}\fi
\noindent\textit{Source reference.} Original reference 04H1: sources/stacks-chow/tags.txt, source lines 5852--5852. Cross-chapter standard prerequisite cited by the author inside the counted span; it is used as a pointer and its own proof is not part of this item.\par\medskip
\ifdefined\hypertarget\hypertarget{sourcepointer-2}{}\fi
\noindent\textit{Source reference.} Original lemma 04IE: sources/stacks-chow/tags.txt, source lines 5900--5900. Cross-chapter standard prerequisite cited by the author inside the counted span; it is used as a pointer and its own proof is not part of this item.\par\medskip
\begin{lemma}
\label{lemma-localize-morphism-ringed-topoi}
Let
$$
f :
(\Sh(\mathcal{C}), \mathcal{O})
\longrightarrow
(\Sh(\mathcal{D}), \mathcal{O}')
$$
be a morphism of ringed topoi. Let $\mathcal{G}$ be a sheaf on $\mathcal{D}$.
Set $\mathcal{F} = f^{-1}\mathcal{G}$.
Then there exists a commutative diagram of ringed topoi
$$
\xymatrix{
(\Sh(\mathcal{C})/\mathcal{F}, \mathcal{O}_\mathcal{F})
\ar[rr]_{(j_\mathcal{F}, j_\mathcal{F}^\sharp)}
\ar[d]_{(f', (f')^\sharp)} & &
(\Sh(\mathcal{C}), \mathcal{O}) \ar[d]^{(f, f^\sharp)} \\
(\Sh(\mathcal{D})/\mathcal{G}, \mathcal{O}'_\mathcal{G})
\ar[rr]^{(j_\mathcal{G}, j_\mathcal{G}^\sharp)} & &
(\Sh(\mathcal{D}), \mathcal{O}')
}
$$
We have $f'_*j_\mathcal{F}^{-1} = j_\mathcal{G}^{-1}f_*$
and $f'_*j_\mathcal{F}^* = j_\mathcal{G}^*f_*$. Moreover, the
morphism $f'$ is characterized by the rule
$$
(f')^{-1}(\mathcal{H} \xrightarrow{\varphi} \mathcal{G})
=
(f^{-1}\mathcal{H} \xrightarrow{f^{-1}\varphi} \mathcal{F}).
$$
\end{lemma}

\begin{proof}
By
Sites, Lemma \ref{sites-lemma-localize-morphism-topoi}
we have the diagram of underlying topoi, the
equality $f'_*j_\mathcal{F}^{-1} = j_\mathcal{G}^{-1}f_*$, and
the description of $(f')^{-1}$.
To define $(f')^\sharp$ we use the map
$$
(f')^\sharp :
\mathcal{O}'_\mathcal{G} =
j_\mathcal{G}^{-1} \mathcal{O}'
\xrightarrow{j_\mathcal{G}^{-1}f^\sharp}
j_\mathcal{G}^{-1} f_*\mathcal{O} =
f'_* j_\mathcal{F}^{-1}\mathcal{O} =
f'_* \mathcal{O}_\mathcal{F}
$$
or equivalently the map
$$
(f')^\sharp :
(f')^{-1}\mathcal{O}'_\mathcal{G} =
(f')^{-1}j_\mathcal{G}^{-1} \mathcal{O}' =
j_\mathcal{F}^{-1}f^{-1}\mathcal{O}'
\xrightarrow{j_\mathcal{F}^{-1}f^\sharp}
j_\mathcal{F}^{-1} \mathcal{O} =
\mathcal{O}_\mathcal{F}.
$$
We omit the verification that these two maps are indeed adjoint
to each other. The second construction of $(f')^\sharp$ shows that
the diagram commutes in the $2$-category of ringed topoi (as the
maps $j_\mathcal{F}^\sharp$ and $j_\mathcal{G}^\sharp$ are identities).
Finally, the equality $f'_*j_\mathcal{F}^* = j_\mathcal{G}^*f_*$
follows from the equality
$f'_*j_\mathcal{F}^{-1} = j_\mathcal{G}^{-1}f_*$
and the fact that pullbacks of sheaves of modules and sheaves of sets agree,
see
Lemma \ref{lemma-localize-ringed-topos}.
\end{proof}

\end{document}
